\chapter{Project Implementation}

\section{Overview of Project Modules}

The repository is organised into two main folders, \texttt{frontend} (the Flutter application) and \texttt{backend} (all services and
pipeline logic), plus the rover firmware, demonstration data and documentation. Table~\ref{tab:modules} maps each module to its
location.

\begin{table}[H]
\centering
\caption{Project modules and their locations.}
\label{tab:modules}
\small
\begin{tabular}{|P{3.2cm}|P{5.8cm}|P{3.8cm}|}
\hline
\textbf{Module} & \textbf{Location} & \textbf{Responsibility}\\\hline
Rover firmware & \texttt{firmware/} & Nano controller, ESP32-CAM server, UART protocol\\\hline
Ingestion & \texttt{backend/gaia/ingestion} & Validation, switches, windows, text chunks\\\hline
Knowledge store & \texttt{backend/gaia/knowledge} & Chunk upserts, hybrid retrieval\\\hline
Pipeline A & \texttt{backend/gaia/pipeline\_a} & Planner, SQL facts, farm RAG, open chat, streaming\\\hline
Pipeline B & \texttt{backend/gaia/pipeline\_b} & VLM, research agent, diagnosis jobs\\\hline
Providers & \texttt{backend/gaia/providers} & LLM and VLM clients, embeddings, web search\\\hline
API and live view & \texttt{backend/gaia/api}, \texttt{core/activity.py} & REST and SSE endpoints, activity log\\\hline
Web interface & \texttt{backend/gaia/web} & Browser version of the app and a retrieval inspector\\\hline
Tools and evaluation & \texttt{backend/scripts} & Rover simulator, bring-up tools, evaluations\\\hline
Mobile application & \texttt{frontend/lib} & Screens, services, models\\\hline
Demonstration data & \texttt{demo\_data/} & Sample history, notes, documents, photos and loader\\\hline
Delivery & \texttt{.github/workflows}, Docker files & APK build, tests, containers\\\hline
\end{tabular}
\end{table}

\section{Data Collection and Preprocessing}

\subsection{From the Sensor to the Backend}

The Nano samples the sensors and sends a checksummed text frame to the ESP32 twice a second. The ESP32 converts the latest frame into
JSON every five seconds, queues it, and posts a batch every 15 seconds (or immediately after a sensor is switched). An example batch
that the ESP32 sends when it has no clock is shown below. Each reading carries only its uptime; the backend reconstructs the time with
Equation~\ref{eq:time}.

\begin{lstlisting}
POST /api/v1/telemetry          X-Device-Key: <key>
{"device_id": "rover-01", "sent_uptime_ms": 130000,
 "sensor_switches": {"dht11": true, "soil_moisture": true, "mq5": true, "mq6": true},
 "readings": [
   {"uptime_ms": 120000,
    "values": {"temperature": 27.5, "humidity": 61.0, "soil_moisture": 43.5,
               "mq5": 210, "mq6": 190, "distance_cm": 91, "battery_v": 7.85},
    "state": {"moving": false, "arm_deployed": true, "soil_valid": true}} ]}
\end{lstlisting}

\subsection{Validation}

For each value the backend normalises the sensor name (for example \texttt{temp} to \texttt{temperature}), rejects non-numeric and
out-of-range values using the sensor registry, drops soil moisture unless the probe state says the arm is deployed and the rover
is still, ignores duplicates (the same device, sensor and time), and rejects timestamps more than five minutes in the future.
The response reports how many values were accepted, rejected and duplicated, with a short list of warnings.

\subsection{Aggregation and Text Generation}

Windows that received new readings are marked ``dirty''. A background job (every 60 seconds) and every chat question run a
\emph{sync}: closed windows are re-aggregated (Equation~\ref{eq:window}), written as text and embedded, and the dirty mark is removed only if no
new reading arrived in the meantime, so late data from the ring buffer is never lost. A window chunk looks like this (a real example
from the demonstration run):

\begin{lstlisting}
Field sensor readings from rover-01 on Saturday 17 October 2026, 16:00-16:15
(afternoon, farm time).
- Air temperature: average 32.8 C (min 32.1, max 33.5, 15 readings) - hot:
  heat stress possible, especially with low soil moisture.
- Relative humidity: average 49% (min 45%, max 52%, 15 readings) - moderate.
- Soil moisture: average 39.2% (min 38.4%, max 40.0%, 10 readings) - adequate.
- MQ5 gas index: average 210 ADC (min 180, max 240, 15 readings) - baseline.
\end{lstlisting}

If a sensor has no statistics in a window, the chunk says why, for example ``MQ5 gas index: not recorded in this window --- sensor
switched off''. Chunks whose text is unchanged are not embedded again (a content hash is compared), and each chunk records the
embedding model that produced its vector, so a change of model never mixes vector spaces (a re-index endpoint re-embeds all chunks).

\subsection{Image Preprocessing}

Uploaded photographs are validated with Pillow (non-images and decompression bombs are rejected), rotated according to their EXIF
orientation, resized to at most 1600 pixels and saved as JPEG. The farmer's optional note is stored as metadata but is
\emph{not} shown to the model, which keeps Pipeline B image-only.

\section{Generative AI Model Selection}

\begin{itemize}
\item \textbf{Chat and planning: Gemini 2.5 Flash.} It supports a response schema (structured output), streaming, and is fast and
      inexpensive enough for several calls per question. It is configurable (\texttt{LLM\_MODEL}); newer models can be
      selected without code changes.
\item \textbf{Vision: Gemini 2.5 Flash.} The same model accepts images, so one API key serves both roles. The vision model is configurable
      separately (\texttt{VLM\_MODEL}).
\item \textbf{Embeddings: gemini-embedding-001 at 768 dimensions}, L2-normalised~\cite{lee2025geminiembedding}. Queries and documents
      use the corresponding task types. An offline hashing embedder is the fallback so that retrieval works without a key.
\item \textbf{Alternatives.} Any OpenAI-compatible service (Groq with Llama models, xAI Grok, local Ollama) can replace the chat or the vision
      model through configuration. The pipelines use only two provider operations, \texttt{complete} (free text) and \texttt{structured}
      (validated JSON, plus a streaming variant), so provider changes are configuration, not code. The local-model path was
      tested end to end with Gemma served by Ollama.
\item \textbf{Web search.} DuckDuckGo through the \texttt{ddgs} package needs no key and is the default; Tavily is used when its key is set.
\end{itemize}

\section{Model Configuration / Fine-Tuning}

No model is fine-tuned. The reasons are those of Table~\ref{tab:approaches}: the knowledge that matters (the farm's own
measurements) changes every few seconds and cannot be absorbed by training, while retrieval gives fresh, citable context. Behaviour is
controlled by prompts, response schemas and a few settings (Table~\ref{tab:settings}).

\begin{table}[H]
\centering
\caption{Main backend settings (\texttt{backend/.env}) and defaults.}
\label{tab:settings}
\small
\begin{tabular}{|P{4.6cm}|P{2.6cm}|P{5.6cm}|}
\hline
\textbf{Setting} & \textbf{Default} & \textbf{Meaning}\\\hline
\texttt{GEMINI\_API\_KEY} & none & Enables chat, vision and embeddings\\\hline
\texttt{WINDOW\_MINUTES} & 15 & Aggregation window (must divide 1440)\\\hline
\texttt{INGEST\_GRACE\_SECONDS} & 90 & Wait for late data before a window is closed\\\hline
\texttt{INGEST\_INTERVAL\_SECONDS} & 60 & Background indexing period\\\hline
\texttt{RAG\_TOP\_K} & 6 & Chunks retrieved per question\\\hline
\texttt{CHAT\_HISTORY\_TURNS} & 6 & Conversation turns sent to the model\\\hline
\texttt{AGENT\_MAX\_ROUNDS} & 2 & Research agent rounds\\\hline
\texttt{AGENT\_QUERIES\_PER\_ROUND} & 3 & Search queries per round\\\hline
\texttt{MAX\_UPLOAD\_MB}, \texttt{VLM\_MAX\_IMAGE\_PX} & 12, 1600 & Photo limits\\\hline
\texttt{DEVICE\_KEY}, \texttt{API\_KEY} & dev key, none & Rover telemetry key; optional app key\\\hline
\texttt{FARM\_TIMEZONE}, \texttt{FARM\_REGION} & Asia/Kolkata, India & Time resolution; region added to web searches\\\hline
\end{tabular}
\end{table}

\noindent\textbf{Prompting rules for answers.} The farm-answer prompt requires that farm-specific statements come from the context and
carry a citation; that \texttt{[F\#]} numbers be quoted and not recomputed; that diagnoses be called possible or suspected; that missing
data be stated plainly with the available data mentioned; that general advice appear only in a final paragraph beginning with
``General guidance:''; that MQ values never be presented as ppm; that every ``unusual reading'' fact be mentioned even if averages look
normal; and that switched-off sensors and disconnected periods be reported exactly without estimating values. The model
returns a structured object (\texttt{answer}, \texttt{cited}, \texttt{used\_general\_knowledge}, \texttt{insufficient\_data},
\texttt{follow\_ups}).

\section{Project Modules}

\subsection{Rover Firmware}

\noindent\textbf{Arduino Nano.} The sketch (about 12.6\,KB of flash, 608\,bytes of RAM) runs motors with a soft ramp (so the gearboxes are not shocked), a 700\,ms
command watchdog, an automatic forward stop at 20\,cm (adjustable), and a soil-arm state machine (retracted, deploying, settling,
deployed, retracting) that lifts the probe before any movement and reports soil only when the arm is deployed, settled and the
rover is still. Sensors are read at their own periods (DHT11 every 2\,s, ultrasonic every 60\,ms). Each sensor has an enable
flag; switching an MQ sensor on restarts its 60-second warm-up, and an optional pin drives an electronic switch in its supply line.

\noindent\textbf{UART protocol.} Frames have the form \texttt{\$BODY*CS} at 57600 baud, where CS is the XOR of the body characters in two
hex digits; frames with a bad checksum are ignored. Commands from the ESP32 are \texttt{D} (drive), \texttt{S} (stop), \texttt{A} (arm),
\texttt{C} (configure), \texttt{W} (sensor switch) and \texttt{P} (ping); the Nano replies \texttt{K} (acknowledged) or \texttt{E} (rejected).
Telemetry is sent at 2\,Hz. Example (MQ-2 not fitted):

\begin{lstlisting}
$T,812,274,610,-1,208,191,435,84,3,7860,56,15,15,15*75
   seq, temp x10, hum x10, mq2, mq5, mq6, soil x10, distance, arm,
   battery mV, flags, installed mask, enabled mask, warm mask
\end{lstlisting}

\noindent\textbf{ESP32-CAM.} It joins the configured Wi-Fi (falling back to its own hotspot), serves the control API on port 80 and the
MJPEG stream on port 81 from a separate server so that streaming never blocks control, and pushes telemetry. Table~\ref{tab:esp32api}
lists the API. The stream is VGA at about 14 frames per second; stills are UXGA (1600$\times$1200). The drive watchdog stops the motors
if no command arrives for 500\,ms, and failed pushes back off exponentially up to two minutes.

\begin{table}[H]
\centering
\caption{HTTP API of the ESP32-CAM.}
\label{tab:esp32api}
\small
\begin{tabular}{|c|P{3.5cm}|P{7.6cm}|}
\hline
\textbf{Method} & \textbf{Path} & \textbf{Purpose}\\\hline
GET & \texttt{:81/stream} & MJPEG live video\\\hline
GET & \texttt{/capture} & One high-resolution JPEG (optional flash)\\\hline
GET & \texttt{/snapshot} & One JPEG at stream resolution\\\hline
GET & \texttt{/api/status} & Wi-Fi, Nano link, drive, arm, battery, push state\\\hline
GET & \texttt{/api/telemetry} & Latest reading in the backend's shape\\\hline
POST & \texttt{/api/drive} & \texttt{left/right}, \texttt{throttle/steer} or \texttt{dir}, repeated every 100--150\,ms\\\hline
POST & \texttt{/api/stop} & Stop at once\\\hline
POST & \texttt{/api/arm} & Deploy or retract the soil probe\\\hline
GET, POST & \texttt{/api/sensors} & Read or set the per-sensor switches (\texttt{\{"all": false\}} is video-only mode)\\\hline
POST & \texttt{/api/config} & Obstacle distance and maximum PWM\\\hline
\end{tabular}
\end{table}

\subsection{Backend: Ingestion and Knowledge Store}

The ingestion service (Algorithm~\ref{alg:ingest}) and the sync job (Algorithm~\ref{alg:sync}) implement the lifecycle of
Figure~\ref{fig:lifecycle}. The knowledge store offers one write operation, \emph{upsert a source}, which replaces all chunks of a logical
source (for example one window or one diagnosis) and skips re-embedding when the content hash is unchanged. Every chunk written
publishes three activity events (created, embedded, stored), which is how the Live view shows all kinds of chunks.

\begin{algorithm}[H]
\caption{Ingest a telemetry batch}
\label{alg:ingest}
\begin{algorithmic}[1]
\Require batch with device id, readings, optional switches, receive time $t^{\mathrm{recv}}$
\State record a heartbeat in the five-minute bucket of $t^{\mathrm{recv}}$
\ForAll{reading $r$ in the batch}
  \State $\hat t \gets$ NTP time if valid, else Equation~\ref{eq:time}, else $t^{\mathrm{recv}}$ for a lone reading
  \If{$\hat t$ is missing or more than 5 minutes in the future} \State reject the values; \textbf{continue} \EndIf
  \ForAll{(name, value) in $r$}
    \State normalise the name; reject non-finite or out-of-range values
    \If{name is soil moisture \textbf{and not} soil valid} \State reject \EndIf
    \State add (sensor, $\hat t$, value) to the row set
  \EndFor
  \If{$r$ carries switches} \State add ($\hat t$, switches) to the switch timeline \EndIf
\EndFor
\State add (batch switches, $t^{\mathrm{recv}}$) to the timeline
\State \textbf{for each} change of a switch state in time order \textbf{do} insert a switch event
\State insert rows ignoring duplicates; mark the affected windows dirty
\State \Return counts of accepted, rejected and duplicate values
\end{algorithmic}
\end{algorithm}

\begin{algorithm}[H]
\caption{Sync: aggregate, write and embed windows}
\label{alg:sync}
\begin{algorithmic}[1]
\Require current time $t$, flag include\_open
\State cutoff $\gets t-\Delta-\text{grace}$
\ForAll{dirty window $w$ in time order}
  \State closed $\gets$ ($w.\mathrm{start}\le$ cutoff)
  \If{\textbf{not} closed \textbf{and not} include\_open} \State \textbf{continue} \EndIf
  \State aggregate the window (Equation~\ref{eq:window})
  \If{\textbf{not} closed} \State \textbf{continue} \Comment{statistics are current, indexing waits}\EndIf
  \State reasons $\gets$ why sensors are missing in $w$ (switch and heartbeat records)
  \State text $\gets$ window chunk from the statistics and reasons
  \State upsert the chunk (embed unless the hash is unchanged)
  \State clear the dirty mark only if it was not re-marked meanwhile
\EndFor
\ForAll{day touched by a closed window}
  \State rebuild the daily summary and the sensor on/off log; upsert both
\EndFor
\end{algorithmic}
\end{algorithm}

\subsection{Pipeline A: Planner, Facts and Retrieval}

\noindent\textbf{Planner.} The planner returns a structured query plan. A rule-based planner with the same output is used when no model
is available or the call fails, and deterministic backstops correct typical model mistakes (for example, a question about a named past period is
never ``current status'').

\noindent\textbf{Facts.} The fact builder computes, from SQL over the window table, the latest readings (with their age and a warning if
stale), the statistics of each requested sensor and period, comparisons and their differences, time-of-day profiles, ``unusual
reading'' facts for short excursions that an average would hide, linked conditions around a diagnosis (plus or minus one hour), log
coverage, and the sensor on/off history described in Section~4.3. If a requested sensor has no data, a ``no data'' fact is followed by
the reason.

\noindent\textbf{Retrieval.} Algorithm~\ref{alg:retrieval} shows the hybrid search. Source types are filtered by intent (only diagnoses
for diagnosis questions) and sensor chunks by time overlap, while diagnoses and notes are searched across all time.

\begin{algorithm}[H]
\caption{Hybrid retrieval with Reciprocal Rank Fusion}
\label{alg:retrieval}
\begin{algorithmic}[1]
\Require query text, $k_{\mathrm{top}}$, optional source types, optional time range
\State $\mathbf q\gets$ normalised embedding of the query (query mode)
\State $V\gets$ top candidates by cosine distance (HNSW index) under the filters
\State $T\gets$ top candidates by PostgreSQL full-text rank under the same filters
\ForAll{chunk $d$ in $V\cup T$}
  \State score$(d)\gets \sum_{r\in\{V,T\}:\,d\in r} 1/(60+\mathrm{rank}_r(d))$
\EndFor
\State \Return the $k_{\mathrm{top}}$ chunks with the highest score, with their vector and lexical ranks
\end{algorithmic}
\end{algorithm}

\subsection{Pipeline A: Grounded Generation and Streaming}

Algorithm~\ref{alg:answer} gives the answer procedure. The streaming decoder (Algorithm~\ref{alg:stream}) extracts the text of the
\texttt{answer} property from the partial JSON the model is writing. It handles escape sequences and Unicode surrogate pairs split
across network chunks. At the end the whole JSON is validated, with one non-streaming regeneration if it is malformed, and the
final answer replaces the streamed draft.

\begin{algorithm}[H]
\caption{Answer a farm question}
\label{alg:answer}
\begin{algorithmic}[1]
\Require question, history, time $t$
\State sync ingestion including the open window
\State plan $\gets$ query plan (LLM or rules); ranges $\gets$ resolve time tokens in farm time
\State facts $\gets$ SQL facts for plan.sensors and ranges, with coverage facts
\State chunks $\gets$ hybrid retrieval of plan.search\_query if plan.search\_knowledge
\If{plan.link\_conditions} \State add facts for $\pm$1 h around each retrieved diagnosis \EndIf
\State context $\gets$ facts $\cup$ chunks, numbered \texttt{[F\#]} and \texttt{[K\#]}
\State out $\gets$ structured LLM answer from context and question (streamed if requested)
\State remove citations whose identifier is not in the context
\State \Return answer, cited sources, grounding flags, follow-up questions
\end{algorithmic}
\end{algorithm}

\begin{algorithm}[H]
\caption{Incremental decoding of the answer field}
\label{alg:stream}
\begin{algorithmic}[1]
\State buffer $\gets$ empty; state $\gets$ seeking
\ForAll{text chunk $c$ produced by the model}
  \State append $c$ to buffer
  \If{state is seeking \textbf{and} buffer contains \texttt{"answer": "}} \State state $\gets$ inside; position $\gets$ end of match \EndIf
  \While{state is inside \textbf{and} a complete character is available at position}
    \State decode one character (handle escapes and surrogate pairs); emit it to the app
    \If{the character is an unescaped quote} \State state $\gets$ done \EndIf
  \EndWhile
\EndFor
\State parse and validate the whole JSON (one non-streaming retry if invalid)
\end{algorithmic}
\end{algorithm}

\subsection{Pipeline B: Vision Analysis and Research Agent}

\noindent\textbf{Vision step.} One call with a fixed system prompt returns JSON with the following fields: whether the image is a plant,
the plant and its part, \texttt{health\_status} (healthy, unhealthy or uncertain), a diagnosis phrased as a possibility, issue
categories (disease, nutrient deficiency, pest damage, fungal, bacterial or viral infection, environmental stress, physical damage,
abnormal coloration or morphology), \emph{evidence} (what is directly visible), \emph{assumptions} (what is inferred), a
differential list, severity, a rough confidence and the image quality. After the call the code re-enforces invariants (a
non-plant image is always ``uncertain'') and attaches a fixed caveat that the confidence is not calibrated.

\noindent\textbf{Research agent.} Algorithm~\ref{alg:agent} shows the loop. The agent is ``one agent'': a bounded loop in code with
the model making the decisions, not a tool-calling framework, so it works identically with every provider.

\begin{algorithm}[H]
\caption{Research agent for a diagnosis}
\label{alg:agent}
\begin{algorithmic}[1]
\Require VLM JSON $a$, search tool, $R_{\max}=2$ rounds, $Q_{\max}=3$ queries per round
\State results $\gets$ empty; asked $\gets$ empty
\For{round $=1$ \textbf{to} $R_{\max}$}
  \State step $\gets$ LLM plan from $a$, the branch policy, asked queries and results so far
  \State queries $\gets$ new queries from step (first $Q_{\max}$); \textbf{if} round $=1$ and none, use fallback queries
  \If{step is done (after round 1) \textbf{or} no queries} \State \textbf{break} \EndIf
  \State run the queries in parallel; keep new URLs only, at most 12 sources; rank institutional sources first
\EndFor
\State findings $\gets$ LLM writer from $a$ and the numbered sources \texttt{[W\#]}
\State remove invalid and invented citations and placeholder rows; \Return findings, sources, trace
\end{algorithmic}
\end{algorithm}

\noindent\textbf{Presenting and writing.} The result page shows the plant, the status, the evidence, the confidence with its caveat,
recapture tips for poor images, the findings and the sources. A ``diagnosis'' chunk and one or more ``agent finding'' chunks are
written with provenance (URLs, domains, model identifiers, whether the answer was web-grounded). Deleting a diagnosis removes
the stored image and its chunks.

\subsection{Sensor Switches and Coverage Engine}

The rover reports its switches; the backend stores only changes in \texttt{sensor\_state\_events} and one row per five-minute bucket in
\texttt{rover\_heartbeats}. The coverage engine (Algorithm~\ref{alg:coverage}) evaluates Equation~\ref{eq:status} for a time range. Consecutive ``on''
buckets (recorded or no reading) are merged for display, because a probe that only measures while parked alternates between the
two statuses all the time, which is normal and not a gap. Descriptions list time spans for ranges up to 36 hours and totals for longer ones.
Rover channels that cannot be switched (distance, battery) are not covered.

\begin{algorithm}[H]
\caption{Coverage of one sensor switch over a range}
\label{alg:coverage}
\begin{algorithmic}[1]
\Require device, switch $s$, range $[a,b)$ clipped to the current time
\State load the switch events before $b$, the buckets with heartbeats, and the buckets with readings (any sensor, and sensor $s$)
\ForAll{five-minute bucket $\beta$ in $[a,b)$}
  \State assign $\mathrm{st}_s(\beta)$ by Equation~\ref{eq:status}
\EndFor
\State merge adjacent buckets with the same status into spans; merge adjacent recorded and no-reading spans
\State \Return spans, and the text: ``readings recorded \ldots; switched off \ldots; rover not connected for \ldots''
\end{algorithmic}
\end{algorithm}

\subsection{Live Activity View}

An in-memory log keeps the last 400 events with increasing identifiers. \texttt{GET /live/state} returns the progress of the open window,
counters, the current switch state and recent events; \texttt{GET /live/stream} streams new events as Server-Sent Events, with a
heartbeat comment every 15 seconds and resume from an identifier. The Live view shows six boxes (rover sensors, server receives,
collecting a window, text chunk, embedding, knowledge store) that light up when an event for them arrives and whose connectors
animate while data flows, a second lane for chat and diagnosis stages, and a feed of recent events.

\subsection{Flutter Mobile Application}

The application has six tabs: \textbf{Overview}, \textbf{Rover} (live MJPEG video, D-pad or joystick with a 120\,ms heartbeat, soil-probe
control, per-sensor switches with a ``video only'' button, and capture), \textbf{Diagnose} (camera, gallery or rover photo, live job
stages, results with sources), \textbf{Gaia AI} (both chat modes, streamed answers, citations that open the source, grounding badges,
follow-up suggestions, a persistent session), \textbf{Sensors} (latest values, charts, an ``off'' badge for switched-off sensors) and
\textbf{Live}. Services (rover link, backend client, repositories) are created once and provided through an inherited widget.
The MJPEG decoder cuts frames at the JPEG start and end markers, which works with any camera boundary format. A small
Server-Sent Events parser handles chunk boundaries, and the chat falls back to the non-streaming endpoint if the server does not
support streaming. Settings hold the backend and rover addresses, an optional API key and the relay switch.

\subsection{Web Interface, Simulator and Tools}

The backend also serves a browser version of the application (overview, rover, diagnosis, chat, sensors) and a knowledge-store
inspector that shows the vector, lexical and fused scores of each chunk. A \emph{rover simulator} implements the same HTTP API as the firmware
(including the sensor switches and heater warm-up), so the whole system can be tested without hardware. \texttt{nano\_console} talks
the UART protocol to a Nano on USB with automated checks, and \texttt{rover\_check} walks the assembled rover's API. A demonstration
loader sends sample data through the real API.

\subsection{Delivery}

A GitHub Actions workflow builds the Android APK from the \texttt{frontend} folder on demand and on changes to it. A second workflow runs
the backend tests (including an embedded PostgreSQL), an offline evaluation, and the application's analysis and tests. Docker Compose starts
the backend with a \texttt{pgvector} PostgreSQL image.

\section{Tools and Technologies Used}

\begin{table}[H]
\centering
\caption{Tools and technologies.}
\label{tab:tools}
\small
\begin{tabular}{|P{3.4cm}|P{9.4cm}|}
\hline
\textbf{Area} & \textbf{Tools}\\\hline
Hardware and firmware & ESP32-CAM, Arduino Nano, Arduino C++, \texttt{arduino-cli}, ArduinoJson, DHT sensor library\\\hline
Backend & Python 3.12, FastAPI, Uvicorn, SQLAlchemy 2, Pydantic 2, httpx, Jinja2\\\hline
Database & PostgreSQL 16, pgvector (HNSW, cosine), full-text search (GIN), embedded server through \texttt{pgserver} for development\\\hline
Generative AI & Gemini API (\texttt{google-genai}), OpenAI-compatible APIs (Groq, Grok, Ollama), DuckDuckGo (\texttt{ddgs}), Tavily\\\hline
Mobile application & Flutter and Dart, \texttt{http}, \texttt{image\_picker}, \texttt{shared\_preferences}, \texttt{url\_launcher}\\\hline
Testing & pytest, \texttt{flutter\_test}, fake provider servers, in-process fake Nano, rover simulator\\\hline
Delivery & GitHub Actions, Docker Compose\\\hline
Documentation & \LaTeX\ (this report), Markdown project documents\\\hline
\end{tabular}
\end{table}

\section{Algorithm Details}

The algorithms of this chapter are collected here with their cost, using $n$ for the readings in a window, $w$ for the windows touched by a
sync, $m$ for the chunks in the store, and $c$ for the number of model calls.

\begin{table}[H]
\centering
\caption{Algorithm summary.}
\label{tab:algsummary}
\small
\begin{tabular}{|P{4.0cm}|P{2.4cm}|P{6.4cm}|}
\hline
\textbf{Algorithm} & \textbf{Cost} & \textbf{Notes}\\\hline
Ingest a batch (Alg.~\ref{alg:ingest}) & $O(r)$ for $r$ values & One bulk insert; duplicates ignored in the database\\\hline
Sync (Alg.~\ref{alg:sync}) & $O(wn)$ plus embedding & One embedding call per changed chunk\\\hline
Hybrid retrieval (Alg.~\ref{alg:retrieval}) & $O(\log m)$ vector search & HNSW index; full-text via GIN index\\\hline
Answer a question (Alg.~\ref{alg:answer}) & 2 model calls & Plan and answer; plus embedding of the query\\\hline
Answer decoding (Alg.~\ref{alg:stream}) & linear in text length & Constant work per character\\\hline
Research agent (Alg.~\ref{alg:agent}) & up to 2 plan calls, 1 writer call & Plus 1 vision call; bounded cost\\\hline
Coverage (Alg.~\ref{alg:coverage}) & $O(b)$ for $b$ buckets & Five-minute resolution\\\hline
\end{tabular}
\end{table}
